% ============================================================================
\clearpage
\section*{v4.9.0-r59 Source-Only: The Correction Strength from a Structural Value}
\addcontentsline{toc}{section}{v4.9.0-r59 Source-Only: The Correction Strength from a Structural Value}
\label{sec:v490-r59-eps-structural}
% ============================================================================

\paragraph{Purpose.}
Layers v4.9.0-r57 and r58 left the corrected trimaximal reading as a fit with two settings.  This layer records the pre-registered test H44 (note
\texttt{llm/Vorab\_Notiz\_H44\_eps\_aus\_Strukturwert.md}, commit \texttt{95566e5}, fixed before computing) of whether a structural value fixes the correction strength
\(\varepsilon\).  To avoid choosing values after the hits of r57 were known, only the candidate list of v4.9.0-r56, fixed before r57, was used.  Earlier layers are
not rewritten.

\noindent\textbf{Source-only status.}
This v4.9.0-r59 layer is source-only.  It does not build a PDF, does not assign a new DOI,
and does not authorize a public release.  The v4.9.0 Version DOI
\href{https://doi.org/10.5281/zenodo.23534954}{10.5281/zenodo.23534954} and the Concept DOI
\href{https://doi.org/10.5281/zenodo.19210728}{10.5281/zenodo.19210728} are unchanged by it.

\subsection*{v4.9.0-r59 results}
\label{subsec:v490-r59-results}

\paragraph{Theoremlet v4.9.0-r59.1 (numerical).}
\begin{enumerate}[leftmargin=2em]
\item[(a)] For none of the fourteen candidates (\(2/9\), \(\beta_{\rm sym}^2\), \(1/3\), \(\pi/9\), \(1/2\), \(\pi/6\), \(\beta_{\rm sym}\), \(1/\sqrt3\), \(2/3\), \(1/\sqrt2\), \(3/4\), \(1\),
\(\sqrt3\), \(\gamma_L\)) does variant A or B of v4.9.0-r57 give a case with \(\theta_{12}\) in the JUNO \(2\sigma\) range and \(\theta_{13}\), \(\theta_{23}\) in the NuFIT
\(3\sigma\) ranges.
\item[(b)] A check after the evaluation: the charged-lepton branches exist at all candidates up to \(\varepsilon=1\); variant A has hits only in two windows,
\(\varepsilon_e=0.03\)--\(0.06\) and \(\varepsilon_e=0.62\)--\(0.64\).
\end{enumerate}

\paragraph{Proof sketch.}
Script \texttt{llm/k99\_test\_h44.py}, output \texttt{llm/k99\_out.txt}; the check in \texttt{llm/k99\_nachkontrolle.py}.  \(\square\)

\paragraph{Verdict and reading.}
H44a is refuted.
\begin{enumerate}[leftmargin=2em]
\item No structural value from the pre-registered list fixes the correction; the reading does not become predictive this way.
\item The data nearly fix \(\varepsilon\) under the reading (\(\approx0.05\) or \(\approx0.63\), each within about \(\pm0.02\)); this is a property of the data, not a structural
value, and searching for a number in these windows now would be numerology.
\item \emph{Correction to v4.9.0-r57.1(a):} the statement ``for \(\varepsilon_e\) from \(0.03\) to \(0.64\)'' gave only the minimum and maximum; the hits lie in the two windows
of (b), with none in between.
\end{enumerate}

\subsection*{v4.9.0-r59 status in the three frames}
\label{subsec:v490-r59-frames}

\begin{center}
\small
\begin{tabular}{@{}>{\raggedright\arraybackslash}p{0.30\textwidth}>{\raggedright\arraybackslash}p{0.62\textwidth}@{}}
\toprule
result & M / T / P\\
\midrule
r59.1(a), (b) & statements on lepton mass matrices in M, the same in T as in v4.9.0-r57; P not evaluated\\
\bottomrule
\end{tabular}
\end{center}

\subsection*{v4.9.0-r59 register entries}
\label{subsec:v490-r59-register}

\begin{center}
\small
\begin{tabular}{@{}>{\raggedright\arraybackslash}p{0.34\textwidth}>{\raggedright\arraybackslash}p{0.58\textwidth}@{}}
\toprule
statement & status after v4.9.0-r59\\
\midrule
\(\varepsilon\) from a structural value & refuted for the pre-registered list (H44)\\
hit windows of v4.9.0-r57 & \(\varepsilon_e\approx0.03\)--\(0.06\) and \(0.62\)--\(0.64\); r57.1(a) corrected\\
\bottomrule
\end{tabular}
\end{center}

\subsection*{v4.9.0-r59 ledger and guard update}
\label{subsec:v490-r59-ledger-guard}

\[
R_{\rm es}^{v490\text{-}r59}=(0,1,0)
\]
for (no structural value for \(\varepsilon\); hit windows located; origin of the correction open).

The layer spends no guard:
\[
R_{\rm glob}^{(v490\text{-}r59,+)}=(0,1,1,1,1),
\qquad
\Gamma_{\rm guard}^{v490\text{-}r59}=(1,1,1,1,1),
\qquad
\beta_4=1/256\ \text{unspent}.
\]
It does not imply exact finite \(C5\) zero closure, exact all-readout terminality,
\(C6/T2\), a global photon-clock/RH theorem, or global all-readout closure.
